% TODO make hw stylesheet
\documentclass[10pt]{article}
\usepackage[letterpaper,top=1in,left=1in,right=1in]{geometry}
\usepackage[dvipsnames,svgnames]{xcolor}
\usepackage[shortlabels]{enumitem}
\usepackage[framemethod=TikZ]{mdframed}
\usepackage{amsmath,amssymb,amsthm}
\usepackage[colorlinks]{hyperref}
\usepackage{multicol}
\usepackage{tabularx}
\usepackage{bbm}

\hypersetup{
    colorlinks=true,
    linkcolor=blue,
    filecolor=magenta,
    urlcolor=magenta,
    pdftitle={MAT 344 HW}
}

\usepackage{mathtools}
\usepackage{thmtools}
\usepackage[textsize=footnotesize]{todonotes}
\usepackage{titling}
\usepackage{cancel}

\reversemarginpar
\mdfdefinestyle{mdbluebox}{roundcorner=10pt,innerbottommargin=9pt,
linecolor=blue,backgroundcolor=TealBlue!5,}
\declaretheoremstyle[headfont=\sffamily\bfseries\color{MidnightBlue},
mdframed={style=mdbluebox},]{thmbluebox}

\mdfdefinestyle{mdbloobox}{frametitlefont=\bfseries,innerbottommargin=8pt,
nobreak=true,backgroundcolor=TealBlue!5,linecolor=blue,}
\declaretheoremstyle[headfont=\bfseries\color{MidnightBlue},
mdframed={style=mdbloobox},headpunct={\\[3pt]},postheadspace=0pt,]{thmbloobox}

\mdfdefinestyle{mdredbox}{frametitlefont=\bfseries,innerbottommargin=8pt,
nobreak=true,backgroundcolor=Salmon!5,linecolor=RawSienna,}
\declaretheoremstyle[headfont=\bfseries\color{RawSienna},
mdframed={style=mdredbox},headpunct={\\[3pt]},postheadspace=0pt,]{thmredbox}

\mdfdefinestyle{mdgreenbox}{linecolor=ForestGreen,backgroundcolor=ForestGreen!5,
linewidth=2pt,rightline=false,leftline=true,topline=false,bottomline=false,}
\declaretheoremstyle[headfont=\bfseries\sffamily\color{ForestGreen!70!black},
mdframed={style=mdgreenbox},headpunct={ --- },]{thmgreenbox}

\mdfdefinestyle{mdorangebox}{linecolor=RedOrange,backgroundcolor=RedOrange!5,
linewidth=2pt,rightline=false,leftline=true,topline=false,bottomline=false,}
\declaretheoremstyle[headfont=\bfseries\sffamily\color{RedOrange!70!black},
mdframed={style=mdorangebox},headpunct={ --- },]{thmorangebox}

\mdfdefinestyle{mdblackbox}{linecolor=black,backgroundcolor=RedViolet!5!gray!5,
linewidth=3pt,rightline=false,leftline=true,topline=false,bottomline=false,}
\declaretheoremstyle[mdframed={style=mdblackbox}]{thmblackbox}

\declaretheoremstyle[spaceabove=6pt,spacebelow=6pt]{boldhead}
\declaretheoremstyle[spaceabove=6pt,spacebelow=6pt,headfont=\normalfont\itshape]{italichead}

\declaretheorem[style=boldhead,name=Problem]{problem}
\declaretheorem[style=boldhead,name=Problem,numbered=no]{problem*}
\declaretheorem[style=boldhead,name=Setup]{setup} % for config geo
\declaretheorem[style=boldhead,name=Example,sibling=problem]{example}
\declaretheorem[style=boldhead,name=Theorem,sibling=problem]{theorem}
\declaretheorem[style=boldhead,name=Corollary,sibling=theorem]{corollary}
\declaretheorem[style=boldhead,name=Proposition,sibling=theorem]{prop}
\declaretheorem[style=boldhead,name=Lemma,numberwithin=problem]{lemma}
\declaretheorem[style=boldhead,name=Theorem,numbered=no]{theorem*}
\declaretheorem[style=boldhead,name=Lemma,numbered=no]{lemma*}
\declaretheorem[style=boldhead,name=Claim,numberwithin=problem]{claim}
\declaretheorem[style=boldhead,name=Claim,numbered=no]{claim*}
\declaretheorem[style=boldhead,name=Remark,numberwithin=problem]{remark}
\declaretheorem[style=boldhead,name=Remark,numbered=no]{remark*}
\declaretheorem[style=boldhead,name=Definition,sibling=theorem]{definition}
\declaretheorem[style=boldhead,name=Definition,numbered=no]{definition*}

\providecommand{\ol}{\overline}
\providecommand{\ceil}[1]{\left\lceil #1 \right\rceil}
\providecommand{\floor}[1]{\left\lfloor #1 \right\rfloor}
\newcommand{\dd}{\,\mathrm{d}}
\providecommand{\ul}{\underline}
\providecommand{\eps}{\varepsilon}
\providecommand{\half}{\frac{1}{2}}
\providecommand{\CC}{\mathbb C}
\providecommand{\FF}{\mathbb F}
\providecommand{\II}{\mathbb I}
\providecommand{\NN}{\mathbb N}
\providecommand{\QQ}{\mathbb Q}
\providecommand{\RR}{\mathbb R}
\providecommand{\ZZ}{\mathbb Z}
\providecommand{\ind}{\mathbbm 1}
\providecommand{\dg}{^\circ}
\providecommand{\ii}{\item}
\providecommand{\setdiff}{\smallsetminus}
\providecommand{\alert}[1]{{\sffamily\textbf{\textcolor{blue}{#1}}}}
\providecommand{\inv}{^{-1}}
\providecommand{\mb}[1]{\mathbf{#1}}
\newcommand{\what}{\widehat}

\newenvironment{soln}{\begin{proof}[Solution]}{\end{proof}}

\providecommand{\pow}{\operatorname{Pow}}
\providecommand{\vocab}[1]{{\textbf{\textcolor{ForestGreen}{#1}}}}
\providecommand{\lcm}{\operatorname{lcm}}
\providecommand{\abs}[1]{\left\lvert #1\right\rvert}
\providecommand{\smabs}[1]{\lvert #1\rvert}
\providecommand{\innerprod}[1]{\langle\!\langle #1\rangle\!\rangle}
\providecommand{\norm}[1]{\left\| #1\right\|}
\providecommand{\smnorm}[1]{\| #1\|}
\providecommand{\gr}{\operatorname{gr}}

\usepackage{mathrsfs}

% place title at top right
\setlength{\droptitle}{-8ex}
\pretitle{\begin{flushright}\Large\bfseries}
\posttitle{\par\end{flushright}}
\preauthor{\begin{flushright}}
\postauthor{\end{flushright}}
\predate{\begin{flushright}}
\postdate{\end{flushright}}

\title{MAT 344 HW01}
\author{\large Aditya Pahuja}
\date{\large Due 02/03}
\begin{document}
\maketitle

\begin{problem}
    [p.6\#2]
    Show that the complex number system can be thought of as
    the field of all polynomials with real coefficients
    modulo the irreducible polynomial $x^2 + 1$.
\end{problem}

\begin{soln}
    Let $\varphi\colon\RR[x]\to\CC$ be the homomorphism determined by
    \begin{equation*}
	\varphi(a) = a\text{ for }a\in\RR,\qquad
	\varphi(x) = i.
    \end{equation*}
    The kernel of this homomorphism contains
    the polynomial $x^2 + 1$ because $\varphi(x^2 + 1) = i^2 + 1 = 0$
    and does not contain any nonzero polynomials of smaller degree because
    $\varphi(a + bx) = a + bi$ is zero if and only if $a=b=0$.
    Since $\RR[x]$ is a principal ideal domain,
    the kernel of $\varphi$ is generated by $x^2 + 1$
    (a nonzero polynomial in the ideal with minimal degree).
    By the first isomorphism theorem for rings,
    $\varphi$ descends to an isomorphism
    between the fields $\RR[x]/(x^2 + 1)$ and $\CC$.
\end{soln}

\begin{problem}
    [p.11\#1]
    Prove that
    \begin{equation*}
	\abs{\frac{a-b}{1-\ol ab}} < 1
    \end{equation*}
    if $\abs a < 1$ and $\abs b < 1$.
\end{problem}

\begin{soln}
    It suffices to prove that $\abs{a-b}^2 < \abs{1-\ol ab}^2$.
    The magnitudes can be rewritten as
    \begin{align*}
        \abs{a-b}^2
	&= (a-b)(\ol a-\ol b) = a\ol a + b\ol b - \ol ab - a\ol b\\
        \abs{1-\ol ab}^2
	&= (1-\ol ab)(1-a\ol b) = 1 + a\ol ab\ol b - \ol ab - a\ol b.
    \end{align*}
    Using $\abs a,\abs b<1$,
    \begin{equation*}
	1 + a\ol ab\ol b - a\ol a-b\ol b
	= (1 - a\ol a)(1 - b\ol b)
	= \underbrace{(1 - \abs a^2)}_{>0}\underbrace{(1 - \abs b^2)}_{>0} > 0
	\implies a\ol a + b\ol b < 1 + a\ol ab\ol b.
    \end{equation*}
    Thus,
    \begin{equation*}
	\abs{a-b}^2
	= a\ol a + b\ol b - \ol ab - a\ol b
	< 1 + a\ol ab\ol b - \ol ab - a\ol b
	= \abs{1-\ol ab}^2
    \end{equation*}
    as desired.
\end{soln}

\begin{problem}
    [p.20\#1]
    Show that $z$ and $z'$ correspond to diametrically opposite points
    on the Riemann sphere if and only if $z\ol{z'} = -1$.
\end{problem}

\begin{soln}
    Let $N$ be the north pole of the Riemann sphere
    and $Z$ and $Z'$ be the points on the sphere
    corresponding to $z$ and $z'$, i.e.
    $Z$ is the second intersection of line $Nz$ with the sphere and
    $Z'$ is the second intersection of line $Nz'$ with the sphere.
    Letting $O$ be the origin of the complex plane (i.e. $0$),
    we have that $\angle NOz = \angle NOz' = 90\dg$.

    If $ZZ'$ is a diameter of the sphere (and $Z$, $Z'$ are not $N$,
    so that neither $z$ nor $z'$ is $\infty$),
    then $\angle zNz' = \angle ZNZ' = 90\dg$
    and $O$ lies on segment $zz'$. The latter statement is true because
    the circle passing through $Z$, $N$, and $Z'$,
    which contains the south pole, gets mapped
    via stereographic projection to a line in the plane containing $z$, $z'$,
    and the image $O$ of the south pole (to be precise,
    $O$ lies on \emph{segment} $zz'$ because the south pole
    lies on the arc connecting $Z$ and $Z'$ that doesn't contain $N$).
    This is all to say that $O$, $N$, $z$, and $z'$ are coplanar. Also,
    $\triangle zON\sim\triangle zNz'$ because both triangles have a right angle
    and share the angle $\angle NzO = \angle Nzz'$,
    and $\triangle z'ON\sim\triangle z'Nz$ by an identical argument.
    This means we have the ratio of distances $zO/ON = NO/Oz'$,
    and since $ON = 1$,
    \[
	\abs z = zO = \frac 1{Oz'} = \frac 1{|z'|}
	\iff |z\ol{z'}| = 1.
    \]
    The fact that $O$ lies on segment $zz'$ additionally means that
    $\frac z{z'} = cz\ol{z'}$ is a negative real number
    (where $c = \frac 1{\abs{z'}^2} > 0$), so $z\ol{z'} = -1$.

    Suppose conversely that $z\ol{z'} = -1$.
    Since this product is a negative real number, $O$ still lies on
    segment $zz'$, meaning $Z$ and $Z'$ are concyclic with
    the north and south poles of the sphere.
    Since the product has magnitude $1$, $Oz/ON = NO/Oz'$,
    which together with $\angle NOz = \angle NOz' = 90\dg$
    implies that $\triangle zON\sim\triangle NOz'$ are similar right triangles.
    Then,
    \begin{equation*}
        \angle zNz' = \angle zNO + \angle ONz'
	= \angle zNO + \angle OzN = 180\dg - \angle NOz = 90\dg,
    \end{equation*}
    where the first equality holds because $O$ is on segment $zz'$,
    so $\angle zNz' = \angle ZNZ' = 90\dg$.
    Together with the fact that $Z$ and $Z'$ lie on a great circle
    of the sphere (the great circle through $N$ and its antipode,
    specifically), we deduce that $ZZ'$ is a diameter of the sphere.
\end{soln}

\begin{problem}
    [p.28\#4]
    Show that an analytic function cannot have a constant absolute value
    without reducing to a constant.
\end{problem}

\begin{soln}
    Let $f = f_1 + if_2$ for real-valued functions $f_1$, $f_2$.
    We are given that $f_1^2 + f_2^2$ is constant, so
    \begin{align*}
	\frac\partial{\partial x}(f_1^2 + f_2^2)
	&= 2f_1\cdot\frac{\partial f_1}{\partial x}
	+ 2f_2\cdot\frac{\partial f_2}{\partial x} = 0\\
	\frac\partial{\partial y}(f_1^2 + f_2^2)
	&= 2f_1\cdot\frac{\partial f_1}{\partial y}
	+ 2f_2\cdot\frac{\partial f_2}{\partial y} = 0.
    \end{align*}

    Let $z\in\CC$ be arbitrary; we will show that
    the above two equations imply $f'(z) = 0$,
    which, since $\CC$ is connected, implies $f$ is constant.
    If both $f_1(z)$ and $f_2(z)$ are zero,
    then $\abs{f(z)} = 0$ everywhere, so $f$ is constantly zero.
    If exactly one of $f_1(z)$ and $f_2(z)$ is zero, say $f_1(z) = 0$,
    then we may divide by $f_2(z)$ to get
    \begin{equation*}
	\frac{\partial f_2}{\partial x}(z)
	= \frac{\partial f_2}{\partial y} = 0.
    \end{equation*}
    By the Cauchy-Riemann equations,
    all the partial derivatives of $f$ are zero, since
    $\frac{\partial f_2}{\partial x} = -\frac{\partial f_1}{\partial y}$
    and $\frac{\partial f_2}{\partial y} = \frac{\partial f_1}{\partial x}$.
    This implies $f'(z) = 0$.

    Finally, suppose that $f_1(z)$ and $f_2(z)$ are both nonzero.
    Rearranging and multiplying the equations yields
    \begin{equation*}
	f_1f_2\cdot
	\frac{\partial f_1}{\partial x}\cdot\frac{\partial f_2}{\partial y} =
	f_1f_2\cdot
	\frac{\partial f_1}{\partial y}\cdot\frac{\partial f_2}{\partial x}
	\implies
	\frac{\partial f_1}{\partial x}\cdot\frac{\partial f_2}{\partial y} =
	\frac{\partial f_1}{\partial y}\cdot\frac{\partial f_2}{\partial x}.
    \end{equation*}
    By the Cauchy-Riemann equations, the left-hand side
    of the latter equation is equal to
    $(\frac{\partial f_1}{\partial x})^2
    = (\frac{\partial f_2}{\partial y})^2 \ge 0$
    and the right-hand side is equal to
    $-(\frac{\partial f_1}{\partial x})^2
    = -(\frac{\partial f_2}{\partial y})^2\le 0$,
    so the LHS and RHS are equal to zero.
    This means that each of the partial derivatives of $f$ is zero,
    hence $f'(z) = 0$.
\end{soln}

\begin{problem}
    [p.28\#5]
    Prove rigorously that the functions $f(z)$ and
    $\ol f(\ol z)$ are simultaneously analytic.
\end{problem}

\begin{soln}
    Let $f\colon U\to\CC$ be an arbitrary function on
    an open subset $U\subseteq\CC$. Then, write
    \begin{equation*}
        f(x+iy) = f_1(x+iy) + if_2(x+iy)
    \end{equation*}
    for functions $f_1,f_2\colon U\to\CC$ and real variables $x,y$.
    Letting $g(z) = \ol f(\ol z) = f_1(x-iy) - if_2(x-iy)$, we have
    \begin{equation*}
	\frac\partial{\partial\ol z}g
	= \half\left(
	\frac{\partial f_1}{\partial x} - \frac{\partial f_2}{\partial y}
        \right)
	- \frac i2\left(
	\frac{\partial f_1}{\partial y} + \frac{\partial f_2}{\partial x}
        \right)
	= \ol{\left(\frac{\partial}{\partial \ol z}f\right)}.
    \end{equation*}
    This means the left-hand side is zero ($\ol f(\ol z)$ is holomorphic)
    if and only if the right-hand side is zero ($f$ is holomorphic).
\end{soln}

\begin{problem}
    [p.28\#6]
    Prove that the functions $u(z)$ and $u(\ol z)$ are
    simultaneously harmonic.
\end{problem}

\begin{soln}
    Let $v(z) = u(\ol z)$ and $u(z) = u_1(z) + iu_2(z)$
    for real-valued functions $u_1$, $u_2$.
    Since conjugation preserves the real part,
    $\frac{\partial^2}{\partial^2 x}v = \frac{\partial^2}{\partial^2 x}u$,
    and
    \begin{equation*}
	\frac{\partial^2}{\partial^2 y}v(z)
	= \frac{\partial}{\partial y}\left(-\frac{\partial u}{\partial y}
	(\ol z)\right) = \frac{\partial^2u}{\partial^2y}(\ol z)
    \end{equation*}
    where the negation in the middle term
    comes from the negation of the imaginary part in $\ol z$.
    Thus $\Delta u = \Delta v$ everywhere, so $u$ is harmonic
    if and only if $v$ is.
\end{soln}










\end{document}
